\section{Case I criterion and the prime predicates}\label{sec:casei}

\begin{theorem}[Cauchy--Genocchi Case I criterion]\label{thm:caseI}
  \lean{CyclotomicNT.caseI_of_not_irregular}
  \leanok
  \uses{def:irrcert, thm:herbrandeigen}
  If $p \nmid B_{p-3}$ (equivalently $p-3$ is not an irregular index), then the
  first case of Fermat's Last Theorem holds for $p$. The load-bearing step is the
  unconditional implication that a Case~I factorization with a good Mirimanoff
  ratio forces $p \mid B_{p-3}$.
\end{theorem}

\begin{definition}[Vandiver prime]\label{def:vandiverprime}
  \lean{IsVandiverPrime}
  \leanok
  \uses{thm:unitindex}
  \code{IsVandiverPrime}~$p$: $p \nmid \hplus$ --- the condition the engine needs,
  characterized via the unit index and certified per prime by the $Q_i$ test.
\end{definition}

\begin{definition}[Regular prime]\label{def:regularprime}
  \lean{IsRegularPrime}
  \leanok
  \uses{thm:notirr, def:vandiverprime}
  \code{IsRegularPrime}~$p$: $p \nmid h(K)$, equivalently $p$ divides no Bernoulli
  numerator $B_2,\dots,B_{p-3}$ (minus part) and $p \nmid \hplus$ (plus part).
\end{definition}
